% Part: exercises
% Chapter: free logics
% Section: tableaus

\documentclass[../../../../include/open-logic-section]{subfiles}

\begin{document}

\olfileid{exe}{frl}{dis}

\olsection{Discussion}

\begin{prob}
Give an English equivalent of each of the following and explain briefly why both
 hold in standard first-order logic.
\begin{enumerate}
\item $\models\exists x(x=x)$
\item $\models\exists(x=a)$
\end{enumerate}

Suppose we extend first-order logic with a necessity operator $\Box$ with the  
necessitation property:
\begin{defn}[Necessitation rule]
If $\Entails !A$, then $\Entails \Box !A$
\end{defn}

\begin{enumerate}
	\item Applying the necessitation rule to the two valid formulas of first-order logic above. which validities do we get? (Give only two.) 
	\item Explain briefly why one would not want to accept the resulting validities.
\end{enumerate}
\end{prob}

\begin{prob}
Negative free logic is committed to primitive predicates.
\begin{enumerate}
	\item Explain what that means.
	\item Explain why that is the case.
	\item Give a reason to think that this is a problem.
\end{enumerate}
\end{prob}

\begin{prob}
Consider the following formula:
$$\forall x (Fx \rightarrow \lfrexists x)$$
\begin{enumerate}
	\item Explain why this is valid in positive and negative free logic.
	\item Give a reason to think that this is a problem.
\end{enumerate}
\end{prob}


\end{document}